\section{The panel-hinge rigidity matrix \texorpdfstring{$R(G,p)$}{R(G,p)}}
\label{sec:molecular-rigidity-matrix}

This chapter builds the panel-hinge (body-hinge) rigidity matrix of
Katoh--Tanigawa~\cite[\S2.2--2.4]{katohTanigawa2011}, the substrate on
which the conjecture's algebraic induction
(\cref{sec:molecular-algebraic-induction}) runs, using the extensor
algebra of \cref{sec:molecular}. Where \cref{sec:body-hinge}
\emph{defined} rigidity by reduction to body-bar frameworks on the
multiplied graph $(\delta-1)\cdot G$ --- a standard-basis witness with no
rank computation, the \emph{existence} form --- this chapter assigns
each edge an explicit $(d-2)$-dimensional affine \emph{hinge} subspace,
builds the rigidity matrix $R(G,p)$ of a realization $p$, and computes
its rank directly. Following \cref{sec:molecular}, the screw space and the
rigidity matrix $R(G,p)$ are built over an arbitrary field $K$;
Katoh--Tanigawa's original argument is the $K = \R$ specialization. The
infinitesimal-motion vectors are $D$-dimensional
\emph{screw centers} $S(v) \in K^{D}$ with $D = \binom{d+1}{2}$,
realized as the degree-$(d-1)$ graded piece $\bigwedge^{d-1}K^{d+1}$ of
the extensor algebra of \cref{sec:molecular-homog}.

\Cref{def:hinge-constraint}, \cref{def:hinge-row-block}, and
\cref{def:rigidity-matrix} build the matrix $R(G,p)$ and its null space
$Z(G,p)$; \cref{lem:trivial-motions-rank-bound} exhibits the
$D$-dimensional space
of trivial (rigid-body) motions, the resulting bound
$\operatorname{rank} R(G,p) \le D(|V|-1)$, and the codimension bound
$D + \operatorname{def}(\tilde G) \le \dim Z(G,p)$ that holds at every
realization, generic or not (the lower-bound half of Tay--Whiteley's
Proposition~1.1, \cref{prop:rigidity-matrix-prop11}).
\Cref{sec:molecular-rigidity-matrix-rank} proves Katoh--Tanigawa's three
rank lemmas --- pinning a body preserves rank (Lemma~5.1,
\cref{lem:rank-delete-vertex}), rank is lower-semicontinuous along a
panel rotation (Lemma~5.2, \cref{lem:rank-rotation-semicont}), and two
hinges of parallel edges in general position force full rank (Lemma~5.3,
\cref{lem:rank-parallel-full}, the base case $|V| = 2$ of the algebraic
induction) --- and \cref{sec:molecular-rigidity-matrix-blocks} proves the
block-rank-addition lemmas underlying the induction's structural cases:
a vertex cut (KT Lemma~6.1, \cref{lem:block-rank-cut}), a shared-body
contraction splice (KT Lemma~6.2,
\cref{lem:rigidityRows-splice-rank-add}), and the pin-a-body splitting
geometry shared by KT Lemmas~6.8 and~6.10
(\cref{lem:rigidityRows-pinned-placement-rank-add}). Equating
$\operatorname{rank} R(G,p) = D(|V|-1)$ with $(\delta,\delta)$-tightness
of $(\delta-1)\cdot G$ (\cref{thm:body-hinge-tay}) needs the analytic
generic-rank theorem of Jackson--Jord\'an
(\cref{prop:rigidity-matrix-prop11}) and is proved alongside the
algebraic induction, once the matroidal bridge
$\operatorname{def}(\tilde G) = \operatorname{corank} M(\tilde G)$
(\cref{sec:molecular-deficiency}) is in hand.

\subsection{Body-hinge frameworks and the rigidity matrix}
\label{sec:molecular-rigidity-matrix-def}

\begin{definition}[Hinge constraint via the supporting extensor]
  \label{def:hinge-constraint}
  \lean{CombinatorialRigidity.Molecular.BodyHingeFramework,
    CombinatorialRigidity.Molecular.BodyHingeFramework.hingeConstraint}
  \leanok
  \uses{def:affine-subspace-extensor}
  A \emph{$d$-dimensional body-hinge framework} $(G,p)$ is a multigraph
  $G = (V,E)$ together with a map $p$ assigning to each edge $e \in E$ a
  $(d-2)$-dimensional affine subspace $p(e)$ of $K^{d}$ --- a
  \emph{hinge}. Identifying an infinitesimal motion of a rigid body with
  a $D$-dimensional screw center, the hinge $p(e)$ supported on the
  $(d-1)$-extensor $C(p(e)) \in \bigwedge^{d-1}K^{d+1} \cong K^{D}$
  (\cref{def:affine-subspace-extensor}) constrains the two screw centers
  $S(u), S(v) \in K^{D}$ of its endpoints $e = uv$ to satisfy
  \[
    S(u) - S(v) \in \operatorname{span} C\bigl(p(e)\bigr).
  \]
\end{definition}

\begin{definition}[Hinge-row block]
  \label{def:hinge-row-block}
  \lean{CombinatorialRigidity.Molecular.BodyHingeFramework.hingeRowBlock,
    CombinatorialRigidity.Molecular.BodyHingeFramework.hingeConstraint_iff_hingeRowBlock,
    CombinatorialRigidity.Molecular.BodyHingeFramework.finrank_hingeRowBlock}
  \leanok
  \uses{def:hinge-constraint}
  For an edge $e$, fix a basis
  $\{r_1(p(e)), \dots, r_{D-1}(p(e))\}$ of the orthogonal complement
  $(\operatorname{span} C(p(e)))^{\perp} \subseteq K^{D}$, and let
  $r(p(e))$ be the $(D-1)\times D$ matrix with these rows. Then a screw
  assignment $S$ meets the hinge constraint
  (\cref{def:hinge-constraint}) at $e = uv$ iff
  $(S(u) - S(v)) \cdot r_i(p(e)) = 0$ for all $1 \le i \le D-1$, by the
  double-annihilator identity
  $(\operatorname{span} C)^{\circ\circ} = \operatorname{span} C$;
  formalized via the dual annihilator
  $(\operatorname{span} C(p(e)))^{\circ} \subseteq (K^{D})^{*}$.
\end{definition}

\begin{definition}[Panel-hinge rigidity matrix]
  \label{def:rigidity-matrix}
  \lean{CombinatorialRigidity.Molecular.BodyHingeFramework.IsInfinitesimalMotion,
    CombinatorialRigidity.Molecular.BodyHingeFramework.infinitesimalMotions,
    CombinatorialRigidity.Molecular.BodyHingeFramework.hingeRow,
    CombinatorialRigidity.Molecular.BodyHingeFramework.rigidityRows,
    CombinatorialRigidity.Molecular.BodyHingeFramework.infinitesimalMotions_eq_dualCoannihilator,
    CombinatorialRigidity.Molecular.BodyHingeFramework.exists_finite_spanning_rigidityRows,
    CombinatorialRigidity.Molecular.BodyHingeFramework.linearIndependent_hingeRow,
    CombinatorialRigidity.Molecular.BodyHingeFramework.exists_independent_rigidityRows_of_edge,
    CombinatorialRigidity.Molecular.BodyHingeFramework.linearIndependent_hingeRow_star,
    CombinatorialRigidity.Molecular.BodyHingeFramework.linearIndependent_hingeRow_forest,
    CombinatorialRigidity.Molecular.BodyHingeFramework.exists_independent_rigidityRows_of_forest,
    CombinatorialRigidity.Molecular.ScrewSpace,
    CombinatorialRigidity.Molecular.screwDim,
    CombinatorialRigidity.Molecular.screwSpace_finrank}
  \leanok
  \uses{def:hinge-row-block}
  The \emph{rigidity matrix} $R(G,p)$ is the $(D-1)|E| \times D|V|$ block
  matrix whose block row for the oriented edge $e = uv$ is $r(p(e))$ in
  the $D$ columns of $u$, $-r(p(e))$ in the $D$ columns of $v$, and zero
  elsewhere. A screw assignment $S \in K^{D|V|}$ is an infinitesimal
  motion of $(G,p)$ iff $S \in \ker R(G,p)$; this null space is written
  $Z(G,p)$, and its dimension is independent of the chosen
  orthogonal-complement basis.
\end{definition}

The formalization keeps the screw space the graded-piece element
$\bigwedge^{d-1}K^{d+1}$ rather than fixing a coordinate matrix: a screw
assignment $S \colon V \to K^{D}$ is an infinitesimal motion when it
meets every per-edge hinge constraint $S(u) - S(v) \in \operatorname{span}
C(p(e))$ --- orientation-independent, since the span is closed under
negation --- and $Z(G,p)$ is the submodule of all such $S$. The block
rows of $R(G,p)$ are the \emph{row functionals} $S \mapsto r(S(u) -
S(v))$, one per oriented edge $e = uv$ and row $r$ of $r(p(e))$; their
span is the row space of $R(G,p)$, and $Z(G,p)$ is exactly the common
kernel of the family, the dual coannihilator of that span.

For an edge $e = uv$ with distinct endpoints, the relative-screw map $S
\mapsto S(u) - S(v)$ is surjective, so the associated row map is
injective: an independent family of hinge-row-block functionals at a
single edge gives an independent family of rigidity rows. Combined with
the $(D-1)$-dimensional hinge-row block of \cref{def:hinge-row-block},
one edge $e = uv$ with $u \ne v$ and $C(p(e)) \ne 0$ therefore
contributes exactly $D - 1$ independent rigidity rows. Independent row
families of \emph{distinct} hinges combine without loss: evaluating a
vanishing combination at a screw supported on a single endpoint isolates
that hinge's block, so per-hinge independent families stay jointly
independent across a star of edges at one body with distinct other
endpoints, and more generally across any spanning forest of hinges
(orienting each hinge from a private ``child'' endpoint: the private
endpoints are distinct, and no hinge's other endpoint coincides with any
hinge's private endpoint). A spanning forest of $|J|$ transversal hinges thus
supplies $(D-1)|J|$ independent rows of $R(G,p)$ --- the count a rigid
subgraph's block-triangular contraction argument later needs to match
against the contraction's inductive rank.

\begin{definition}[Degree of freedom; generic realization]
  \label{def:dof-generic}
  \lean{CombinatorialRigidity.Molecular.BodyHingeFramework.IsInfinitesimallyRigid}
  \uses{lem:trivial-motions-rank-bound}
  An infinitesimal motion $S$ is \emph{trivial} when $S(u) = S(v)$ for
  all $u, v$; $(G,p)$ is \emph{infinitesimally rigid} when all
  infinitesimal motions are trivial, equivalently
  $\operatorname{rank} R(G,p) = D(|V|-1)$. The \emph{degree of freedom}
  of $(G,p)$ is $D(|V|-1) - \operatorname{rank} R(G,p)$. A realization is
  \emph{generic} when $R(G,p)$ and all its edge-induced submatrices
  attain their maximum ranks over realizations of $G$. Formalized
  basis-free as the inclusion $Z(G,p) \subseteq \{\text{trivial motions}\}$
  (\texttt{IsInfinitesimallyRigid}), the reverse of the always-true
  $\{\text{trivial}\} \subseteq Z(G,p)$ of
  \cref{lem:trivial-motions-rank-bound}.
\end{definition}

\begin{lemma}[The $D$ trivial motions and the rank bound]
  \label{lem:trivial-motions-rank-bound}
  \lean{CombinatorialRigidity.Molecular.BodyHingeFramework.IsTrivialMotion,
    CombinatorialRigidity.Molecular.BodyHingeFramework.isInfinitesimalMotion_of_isTrivialMotion,
    CombinatorialRigidity.Molecular.BodyHingeFramework.trivialMotions,
    CombinatorialRigidity.Molecular.BodyHingeFramework.trivialMotions_le_infinitesimalMotions,
    CombinatorialRigidity.Molecular.BodyHingeFramework.trivialMotions_eq_range_const,
    CombinatorialRigidity.Molecular.BodyHingeFramework.infinitesimalMotions_eq_trivialMotions_iff,
    CombinatorialRigidity.Molecular.BodyHingeFramework.finrank_trivialMotions,
    CombinatorialRigidity.Molecular.BodyHingeFramework.finrank_screwAssignment,
    CombinatorialRigidity.Molecular.BodyHingeFramework.screwDim_add_deficiency_le_finrank_infinitesimalMotions}
  \leanok
  \uses{def:rigidity-matrix, def:D-deficiency}
  For $1 \le i \le D$ let $S^*_i$ be the constant screw assignment with
  $i$-th coordinate $1$ at every vertex. Each $S^*_i$ lies in $Z(G,p)$
  and is trivial, and $\{S^*_1, \dots, S^*_D\}$ is linearly independent
  and spans the space of trivial infinitesimal motions. Hence
  $\operatorname{rank} R(G,p) \le D|V| - D = D(|V|-1)$, with equality iff
  $(G,p)$ is infinitesimally rigid. More is true, and holds at every
  realization regardless of genericity: a partition $P$ of $V$
  contributes its full part-constant freedoms, giving the codimension
  bound $D + \operatorname{def}(\tilde G) \le \dim Z(G,p)$ --- the
  lower-bound half of Proposition~1.1
  (\cref{prop:rigidity-matrix-prop11}).
\end{lemma}
\begin{proof}
  \leanok
  A constant assignment has $S(u) - S(v) = 0$, which lies in every
  hinge constraint, so $S^*_i \in Z(G,p)$ and is trivial, giving
  $\{\text{trivial}\} \subseteq Z(G,p)$; the
  $\{S^*_i\}$ are linearly independent across the $D$ coordinates and a
  trivial motion is determined by its single common value, so they span
  the trivial motions, which are exactly the constant assignments
  (\texttt{trivialMotions}, the diagonal of $V \to K^{D}$). The
  constant-assignment map is thus a linear isomorphism
  $K^{D} \cong \{\text{trivial}\}$ for nonempty $V$, so the trivial
  kernel has dimension exactly $D$ inside the $D|V|$-dimensional
  screw-assignment space. The $D$-dimensional trivial kernel therefore
  forces $\operatorname{rank} R(G,p) \le D|V| - D = D(|V|-1)$.

  For the codimension bound, fix a labeling $f$ whose fibers are the
  parts of $P$, and consider the part-constant assignments $W_f$ (one
  screw center per part), of dimension $D \cdot |P|$. Each of the
  $d_G(P)$ edges crossing $P$ imposes one relative-screw constraint
  $S(u) - S(v) \in \operatorname{span} C(e)$, cutting the part-constant
  motions $W_f \cap Z(G,p)$ down by at most $D - 1$ per crossing edge (the
  hinge is nondegenerate, $C(e) \neq 0$); non-crossing edges impose
  nothing on a part-constant assignment. Rank-nullity over $W_f$ thus
  gives $\dim (W_f \cap Z) \ge D|P| - (D-1)\,d_G(P) = D + \operatorname{def}_{\tilde
  G}(P)$. Maximizing over $P$ (a finite supremum) attains
  $\operatorname{def}(\tilde G)$, and $W_f \cap Z \subseteq Z$ transfers the
  bound to $\dim Z(G,p) \ge D + \operatorname{def}(\tilde G)$.
\end{proof}

\subsection{Rank lemmas}
\label{sec:molecular-rigidity-matrix-rank}

\begin{lemma}[Deleting one body's columns preserves rank; KT Lemma 5.1]
  \label{lem:rank-delete-vertex}
  \lean{CombinatorialRigidity.Molecular.BodyHingeFramework.trivialMotions_sup_pinnedMotions,
    CombinatorialRigidity.Molecular.BodyHingeFramework.trivialMotions_inf_pinnedMotions_eq_bot,
    CombinatorialRigidity.Molecular.BodyHingeFramework.finrank_pinnedMotions_add_screwDim}
  \leanok
  \uses{def:rigidity-matrix, lem:trivial-motions-rank-bound}
  Pinning one body (deleting the $D$ columns of a single vertex $v$ from
  $R(G,p)$) does not change the rank:
  $\operatorname{rank} R(G,p) = \operatorname{rank} R(G,p; V \setminus v)$.
  Consequently rank computations may always fix one body, the most
  reused reduction in the algebraic induction.
\end{lemma}
\begin{proof}
  \leanok
  The $D$ trivial motions (\cref{lem:trivial-motions-rank-bound}) let
  any motion be normalized to vanish on the pinned body, so the column
  space is unchanged; this is the ``pinning a body preserves the motion
  space'' fact of White--Whiteley~\cite{whiteWhiteley1987}. Concretely,
  the null space $Z(G,p)$ is the internal direct sum of the trivial
  motions and the body-$v$-pinned motions (those vanishing on $v$): any
  motion $S$ splits as the constant $u \mapsto S(v)$ plus the pinned
  remainder $u \mapsto S(u) - S(v)$, which is again a motion because the
  hinge constraint only sees the relative screws $S(u) - S(w)$
  (\texttt{isInfinitesimalMotion\_sub\_const}); the two summands meet only
  at $0$. Hence $\operatorname{finrank}(\text{pinned at } v) + D =
  \operatorname{finrank} Z(G,p)$, the codimension form of
  rank-preservation.
\end{proof}

\begin{lemma}[Two hinges of parallel edges give the full block; KT Lemma 5.3]
  \label{lem:rank-parallel-full}
  \lean{CombinatorialRigidity.Molecular.BodyHingeFramework.span_inf_span_eq_bot_of_linearIndependent,
    CombinatorialRigidity.Molecular.BodyHingeFramework.eq_of_hingeConstraint_two_parallel}
  \leanok
  \uses{def:rigidity-matrix, lem:extensor-independence}
  If two edges of $G$ join the same pair of bodies with hinges in
  \emph{general position}, the two $(D-1) \times D$ blocks together have
  rank $D$ --- the constraint forces the relative screw to zero. The
  base case $|V| = 2$ of the conjecture's induction is this lemma.
\end{lemma}
\begin{proof}
  \leanok
  Each hinge constrains the relative screw to the span of its supporting
  extensor (\cref{def:rigidity-matrix}); two extensors in general
  position are linearly independent, by the extensor-independence engine
  (\cref{lem:extensor-independence}, specialized to two hinges), so their
  one-dimensional constraint spans meet only at $0$
  (\texttt{span\_inf\_span\_eq\_bot\_of\_linearIndependent}). Hence a
  screw satisfying both constraints has zero relative part, the combined
  block has the full rank $D$.
\end{proof}

\begin{lemma}[Lower-semicontinuity of rank under a panel rotation; KT Lemma 5.2]
  \label{lem:rank-rotation-semicont}
  \lean{CombinatorialRigidity.Molecular.BodyHingeFramework.infinitesimalMotions_mono_of_span_le,
    CombinatorialRigidity.Molecular.BodyHingeFramework.finrank_infinitesimalMotions_le_of_span_le}
  \leanok
  \uses{def:rigidity-matrix, def:dof-generic}
  Along a one-parameter family of realizations rotating a single panel,
  the rank of $R(G,p)$ is lower-semicontinuous: it cannot drop at a
  generic parameter. Equivalently, if some realization in the family
  attains rank $\rho$, so does the generic one. In the basis-free
  monotonicity form: if two realizations $p, p'$ of the same multigraph
  satisfy $\operatorname{span} C(p(e)) \subseteq \operatorname{span} C(p'(e))$
  at every edge --- $p'$ the more general (generic) realization, with the
  larger supporting spans --- then $Z(G,p') \subseteq Z(G,p)$ and hence
  $\operatorname{finrank} Z(G,p') \le \operatorname{finrank} Z(G,p)$, i.e.\
  $\operatorname{rank} R(G,p) \le \operatorname{rank} R(G,p')$.
\end{lemma}
\begin{proof}
  \leanok
  The entries of $R(G,p)$ are polynomials in the panel coordinates, so a
  nonvanishing $\rho \times \rho$ minor at one parameter is a nonzero
  polynomial and hence nonvanishing at a generic parameter --- the
  genericity argument is run directly, rather than through an analytic
  perturbation. The formalization proves the basis-free monotonicity
  statement: a refinement of the constraint spans ($\operatorname{span}
  C(p(e)) \subseteq \operatorname{span} C(p'(e))$ for every edge) makes
  every motion of the more-constrained $p'$ a motion of $p$, so
  $\operatorname{finrank} Z$ only shrinks under refinement; the generic
  member of the rotation family has the largest supporting spans and
  thus the maximal rank.
\end{proof}

\subsection{Block-rank-addition lemmas for the case analyses}
\label{sec:molecular-rigidity-matrix-blocks}

The algebraic induction's structural cases each reduce a rank count on
$G$ to rank counts on smaller pieces plus a block correction.
\Cref{lem:block-rank-cut} is the vertex-cut form (KT Lemma~6.1); the
remaining lemmas cover the two body-sharing geometries a vertex cut
cannot express --- \emph{collapse}, where a rigid subgraph $H$ and its
contraction $G/E(H)$ share one body (KT Lemma~6.2,
\cref{lem:rigidityRows-splice-rank-add}), and \emph{pin-a-body}, where a
new block of rows and an old block share a body's screw column (KT
Lemmas~6.8 and~6.10, \cref{lem:rigidityRows-pinned-placement-rank-add}).
The supporting lemmas discharge the two geometries' respective
hypotheses: injectivity of the collapse's column-deletion map
(\cref{lem:extProj-preserves-rank-of-inter}), rank invariance under equal
motion spaces (\cref{lem:span-rigidityRows-eq-of-motions-eq}), and the
deficiency-aware rank transport a collapsed leg needs when its
contraction is not yet fully rigid
(\cref{lem:rank-transport-relabel-of-le-finrank} and
\cref{lem:rank-polynomial-IH-relabel-proj}).

\begin{lemma}[Block-rank addition for a vertex cut; KT Lemma 6.1]
  \label{lem:block-rank-cut}
  \lean{CombinatorialRigidity.Molecular.BodyHingeFramework.le_finrank_span_rigidityRows_of_cut}
  \leanok
  \uses{def:rigidity-matrix, def:hinge-row-block}
  Let $(G,p)$ be a body-hinge framework with all nonzero supporting extensors,
  and let $V_1 \subsetneq V(G)$ be a nonempty proper vertex subset with cut
  $C = \operatorname{cutEdges}(G,V_1)$ satisfying $|C| \le 1$.  Write
  $V_2 = V(G) \setminus V_1$ and let $F_i = (G[V_i],p)$ be the induced
  subframeworks.  Then
  \[
    \operatorname{rank} R(G[V_1],p) + (D-1)\,|C|
    + \operatorname{rank} R(G[V_2],p)
    \;\le\;
    \operatorname{rank} R(G,p).
  \]
  This is the vertex-disjoint block-rank-addition lower bound for a
  cut-edge decomposition; it replaces the isometry argument of
  Katoh--Tanigawa's Lemma~6.1~\cite{katohTanigawa2011} with an argument
  via three mutually disjoint subspaces of the rigidity-row span at a
  fixed placement.
\end{lemma}
\begin{proof}
  \leanok
  \uses{def:rigidity-matrix, def:hinge-row-block}
  When $|C| = 0$ the two sides are disconnected: $S_1, S_2 \le S$ and
  $S_1 \cap S_2 = \bot$ (a flow-sum argument separates them), giving
  $\operatorname{finrank}(S_1) + \operatorname{finrank}(S_2)
  \le \operatorname{finrank}(S)$.
  When $|C| = 1$ the single cut edge $e_c$ contributes a $(D-1)$-dimensional
  subspace $S_c$ (from the hinge-row block of $e_c$, nonzero by hypothesis);
  the three subspaces $S_1, S_c, S_2$ are pairwise disjoint, and a
  calc chain of disjoint-rank additions closes the inequality.
\end{proof}

\begin{lemma}[Block-triangular rank-addition for a shared-body splice; KT Lemma 6.2]
  \label{lem:rigidityRows-splice-rank-add}
  \lean{CombinatorialRigidity.Molecular.BodyHingeFramework.le_finrank_span_rigidityRows_of_splice}
  \leanok
  \uses{def:rigidity-matrix, def:hinge-row-block}
  Let $F$, $F_H$, $F_c$ be body-hinge frameworks on the same bodies and hinges.
  Let $S = \operatorname{span}(R(F))$, $S_H = \operatorname{span}(R(F_H))$,
  $S_c = \operatorname{span}(R(F_c))$ be the corresponding rigidity-row spans.
  Let $D$ be a linear endomorphism of the row dual space (in practice, the
  dual of the projection that collapses the contracted body; see
  \cref{sec:molecular-algebraic-induction-caseI}). Suppose:
  \begin{enumerate}
    \item $S_H \le S$ and $S_H \le \ker D$ (the hinge block lies in $S$ and is killed by $D$);
    \item $D(S_c) \le D(S)$ (the contraction block maps into the image of $S$);
    \item $D$ is injective on $S_c$ (so $\operatorname{rank}(S_c) = \operatorname{rank}(D(S_c))$).
  \end{enumerate}
  Then
  \[
    \operatorname{rank}(S_H) + \operatorname{rank}(S_c) \;\le\; \operatorname{rank}(S).
  \]
\end{lemma}
\begin{proof}
  \leanok
  \uses{def:rigidity-matrix, def:hinge-row-block}
  Rank-nullity applied to $D$ restricted to $S$ gives
  $\operatorname{rank}(D(S)) + \operatorname{rank}(S \cap \ker D) = \operatorname{rank}(S)$.
  By hypothesis (1), $S_H \le S \cap \ker D$, so
  $\operatorname{rank}(S_H) \le \operatorname{rank}(S \cap \ker D)$.
  By hypotheses (2) and (3),
  $\operatorname{rank}(S_c) = \operatorname{rank}(D(S_c)) \le \operatorname{rank}(D(S))$.
  Adding the two inequalities and using rank-nullity gives the bound.
\end{proof}

\begin{lemma}[Span-level pinned-placement block-rank lower bound; KT Lemmas 6.8 and 6.10]
  \label{lem:rigidityRows-pinned-placement-rank-add}
  \lean{CombinatorialRigidity.Molecular.BodyHingeFramework.le_finrank_span_rigidityRows_of_pinned_placement,
    CombinatorialRigidity.Molecular.BodyHingeFramework.le_finrank_span_rigidityRows_of_pinned_placement_augment}
  \leanok
  \uses{def:rigidity-matrix, lem:case-II-placement-block-independent, lem:case-III-conditional-block}
  The span-transport analogue of \cref{lem:rigidityRows-splice-rank-add}
  for the \emph{pin-a-body} (vertex-splitting) geometry, as opposed to
  that lemma's \emph{collapse} (projected-column) geometry --- Katoh--Tanigawa's
  own Lemma~6.2 (contraction) versus Lemma~6.8 (splitting) distinction.
  Let $F$ be a body-hinge framework, $v$ a body, and write $S =
  \operatorname{span}(R(F))$ for its rigidity-row span. Let $\{r_n^i\}_{i \in \iota_n}$
  be a \emph{new block} of row functionals and $\{r_o^j\}_{j \in \iota_o}$ an \emph{old
  block}, with:
  \begin{enumerate}
    \item $r_n$ is independent when restricted to screw assignments
      supported at $v$ alone;
    \item $r_o$ vanishes on every screw assignment supported at $v$
      alone;
    \item $r_o$ is independent;
    \item every $r_n^i \in S$ and every $r_o^j \in S$.
  \end{enumerate}
  Then $|\iota_n| + |\iota_o| \le \operatorname{rank}(S)$. The $+1$ augment variant adds one
  further candidate row $w \in S$ to the new block, independent through $v$'s column alongside
  $r_n$, and concludes $|\iota_n| + 1 + |\iota_o| \le \operatorname{rank}(S)$ --- the extra
  row supplying the lift from $D(|V|-1)-1$ to the full $D(|V|-1)$ that
  Case~III needs (KT eq.~(6.29)).
\end{lemma}
\begin{proof}
  \leanok
  \uses{lem:case-II-placement-block-independent, lem:case-III-conditional-block}
  The block-triangular pin-a-body split
  (\cref{lem:case-II-placement-block-independent}; for the augment,
  \cref{lem:case-III-conditional-block}) makes the family obtained by
  concatenating $r_n$ and $r_o$ linearly independent. Its span lies in
  $S$ by hypothesis~(4), so its dimension equals the block-size sum and
  is bounded by $\operatorname{rank}(S)$. This is the rank skeleton
  shared by KT Lemma~6.8 and Lemma~6.10's arguments, abstracted away
  from any specific construction: no new linear algebra is needed once
  hypothesis~(4) is discharged, which the callers do differently (the
  split-off $e_0 = e_a + e_b$ row decomposition for Lemma~6.8, the
  relabelled hinge-row span for Case~III). Unlike a literal placement
  bound, the conclusion keys on \emph{membership in} $S$ rather than
  literal rigidity-row membership, so it applies to every reduction
  move --- collapse, split-off, or relabel --- whose rows land only in
  $S$. See \cite[\S6.3]{katohTanigawa2011}.
\end{proof}

\begin{lemma}[The exterior-column projection preserves the rigidity-row rank at any rank]
  \label{lem:extProj-preserves-rank-of-inter}
  \lean{CombinatorialRigidity.Molecular.infinitesimalMotions_sup_range_extProj_eq_top_of_inter_eq_singleton,
    CombinatorialRigidity.Molecular.BodyHingeFramework.injOn_extProj_dualMap_rigidityRows_of_inter_eq_singleton,
    CombinatorialRigidity.Molecular.BodyHingeFramework.finrank_span_rigidityRows_map_extProj_dualMap_of_inter_eq_singleton}
  \leanok
  \uses{def:rigidity-matrix, lem:rank-delete-vertex, lem:claim-6-4}
  Let $F$ be a body-hinge framework and $t$ a body set meeting the vertex
  set in exactly one body, $V(F.\mathrm{graph}) \cap t = \{r\}$.  Let
  $D = (\operatorname{extProj} t)^{\!*}$ be the dual of the projection that
  zeroes the $t$-columns.  Then $D$ is injective on the rigidity-row span
  $S = \operatorname{span}(R(F))$, so the projection preserves its rank,
  $\operatorname{rank}(S) = \operatorname{rank}(D(S))$.  This is the
  injectivity hypothesis (3) of \cref{lem:rigidityRows-splice-rank-add},
  discharged at a \emph{deficient} contraction leg where the
  full-rigidity-gated form (\cref{lem:claim-6-4}) is unavailable.
\end{lemma}
\begin{proof}
  \leanok
  \uses{def:rigidity-matrix, lem:rank-delete-vertex}
  Injectivity on $S$ is $S \cap \ker D = \bot$, equivalently
  $Z \sqcup W = \top$ for $Z$ the infinitesimal motions and
  $W = \operatorname{range}(\operatorname{extProj} t)$ the screw
  assignments vanishing on $t$ (by the dual-annihilator calculus:
  $\ker D = W^\circ$, $S = Z^\circ$, so $S \cap \ker D = (Z \sqcup W)^\circ$).
  Unlike the rigid form, $Z \sqcup W = \top$ holds at \emph{any} rank by an
  explicit decomposition needing no rank count: any assignment $\sigma$
  splits as $z + (\sigma - z)$ with $z(a) = \sigma(r)$ for $a \in V(G)$ and
  $z(a) = \sigma(a)$ otherwise; $z$ is a motion (every constraint has both
  endpoints in $V(G)$, where $z$ is the constant $\sigma(r)$, so the
  relative screw is $0$), and $\sigma - z$ vanishes on $t = \{r\} \cup
  (t \setminus V(G))$ since $V(G) \cap t = \{r\}$.  This is the algebraic
  content KT~\cite{katohTanigawa2011} invokes through Lemma~5.1
  (\cref{lem:rank-delete-vertex}): deleting the single shared body $r$'s
  columns preserves rank, the only motion lost to the projection being the
  trivial $\sigma(r)$-shift recovered by $z$.  Rank-nullity for $D$
  restricted to $S$ then gives $\operatorname{rank}(D(S)) = \operatorname{rank}(S)$.
\end{proof}

\begin{lemma}[Rigidity-row span is the motions' annihilator; equal motions give equal rank]
  \label{lem:span-rigidityRows-eq-of-motions-eq}
  \lean{CombinatorialRigidity.Molecular.BodyHingeFramework.span_rigidityRows_eq_dualAnnihilator_infinitesimalMotions,
    CombinatorialRigidity.Molecular.BodyHingeFramework.span_rigidityRows_eq_of_infinitesimalMotions_eq}
  \leanok
  \uses{def:rigidity-matrix}
  Over a finite body set, the rigidity-row span of a body-hinge
  framework is the dual annihilator of its infinitesimal-motion space,
  $\operatorname{span}(R(F)) = Z(F)^\circ$.  Consequently two frameworks
  with the same motion space have the same rigidity-row span --- hence
  the same rank --- at any deficiency, with no rigidity hypothesis.
\end{lemma}
\begin{proof}
  \leanok
  \uses{def:rigidity-matrix}
  The motion space is the dual coannihilator of the row span
  (\texttt{infinitesimalMotions\_eq\_dualCoannihilator}); over a finite
  body set the screw-assignment dual is finite-dimensional, so the
  double-annihilator identity closes
  $\operatorname{span}(R(F)) = Z(F)^\circ$.  Equal motion spaces then
  give equal annihilators, hence equal spans and equal finrank.
\end{proof}

\begin{lemma}[Deficiency-aware rank transport across a collapse relabel]
  \label{lem:rank-transport-relabel-of-le-finrank}
  \lean{CombinatorialRigidity.Molecular.PanelHingeFramework.finrank_span_rigidityRows_ofNormals_relabel_eq}
  \leanok
  \uses{def:rigidity-matrix, lem:span-rigidityRows-eq-of-motions-eq, def:genuine-hinge-realization}
  Let $H \le G$ be a subgraph with representative body $r \in V(H)$, let
  $G_c := G \setminus E(H)$ be the surviving-edge subgraph, and let $f
  \colon V(G) \to V(G)$ collapse every body of $V(H)$ onto $r$ and fix
  every other body, so that $f(G_c) = G/E(H)$ is the contracted graph.
  Given the contraction's inductive realization --- a general-position
  panel framework on $f(G_c)$ attaining the rigidity-row rank
  $D(|V(f(G_c))|-1) - \operatorname{def}(f(G_c))$, at any deficiency ---
  and a parent selector recording $G_c$'s links, the free-normal panel
  framework on $f(G_c)$ at the \emph{relabel selector} $e \mapsto
  (f(\text{ends}(e)_1), f(\text{ends}(e)_2))$ is again in general
  position and attains \emph{the same} rigidity-row rank. This is the
  deficiency-tolerant sibling of the rigid contraction-leg transport
  (\cref{lem:claim-6-4}): it transports the contraction's
  possibly-deficient rank ($\operatorname{def} = k > 0$) across the
  relabel, where the full-rigidity argument does not apply.
\end{lemma}
\begin{proof}
  \leanok
  \uses{lem:span-rigidityRows-eq-of-motions-eq}
  The relabel selector and the inductive realization's own selector
  record the same unordered link of every edge of $f(G_c)$, so they
  agree up to a swap of endpoints; since the hinge constraint is
  membership in the (sign-insensitive) line through the supporting
  extensor, swapping endpoints fixes every constraint, hence the whole
  motion space.  Equal motion spaces give equal rigidity-row rank
  (\cref{lem:span-rigidityRows-eq-of-motions-eq}); general position is a
  property of the normals alone, unchanged by the selector.
\end{proof}

\begin{lemma}[Deficiency-aware rank polynomial for the relabeled contraction leg (superseded)]
  \label{lem:rank-polynomial-IH-relabel}
  \lean{CombinatorialRigidity.Molecular.PanelHingeFramework.exists_rankPolynomial_of_IH_relabel_linking}
  \leanok
  \uses{lem:rank-transport-relabel-of-le-finrank, lem:rank-polynomial-of-le-finrank}
  \emph{Superseded; no live consumer.} This lemma produces the
  contraction framework's \emph{full-span} surviving rank --- the input
  \cref{lem:rigidityRows-splice-rank-add} was designed to consume. The
  realization this was built for cannot use that lemma directly, because
  its surviving-edge support extensors are not in general position for a
  free-normal panel framework; the working argument instead reads the
  \emph{exterior-projected} surviving rows
  (\cref{lem:rank-polynomial-IH-relabel-proj}). The statement below is
  true and its proof's shared step (\cref{lem:rank-transport-relabel-of-le-finrank})
  remains in use; the lemma itself is retained for the record.

  Under the setup of \cref{lem:rank-transport-relabel-of-le-finrank} (with $f(G_c)$ loopless),
  there exists a natural number $N$ with
  $(N : \mathbb{Z}) = D(|V(f(G_c))| - 1) - \operatorname{def}(f(G_c))$
  and a nonzero polynomial $Q$ over the body-normal parameters of the original graph
  such that at every $Q$-non-root seed $q$, the rigidity-row span of
  $\operatorname{ofNormals}(f(G_c), e \mapsto (f(\text{ends}(e)_1), f(\text{ends}(e)_2)), q)$
  has rank at least $N$.
\end{lemma}
\begin{proof}
  \leanok
  \uses{lem:rank-transport-relabel-of-le-finrank, lem:rank-polynomial-of-le-finrank}
  Apply \cref{lem:rank-transport-relabel-of-le-finrank} to obtain a
  witness seed $q_0$ at which the rigidity-row rank equals $N$ exactly
  and the framework is in general position. Then apply the
  deficiency-aware rank polynomial (\cref{lem:rank-polynomial-of-le-finrank})
  at $q_0$ with the trivial bound $N \le N$: the support extensor is
  nonzero at every link of $f(G_c)$ because general position forces it
  and looplessness ensures distinct endpoints. The resulting polynomial
  $Q$ is nonzero (its evaluation at $q_0$ is nonzero), and
  it transfers the rank lower bound $N$ to every non-root seed.
\end{proof}

\begin{lemma}[Deficiency-aware exterior-projected rank polynomial for the surviving block]
  \label{lem:rank-polynomial-IH-relabel-proj}
  \lean{CombinatorialRigidity.Molecular.PanelHingeFramework.exists_rankPolynomial_of_IH_relabel_linking_set_proj}
  \leanok
  \uses{lem:rank-transport-relabel-of-le-finrank, lem:extProj-preserves-rank-of-inter, lem:rank-polynomial-of-le-finrank}
  Let $H \le G$ be a proper rigid subgraph with representative body
  $r \in V(H) \subseteq V(G)$, and let the contraction $G/E(H)$ be a minimal
  $k'$-dof-graph carrying its inductive general-position realization (the
  deficiency $k' > 0$ permitted).  Write $D = (\operatorname{extProj} V(H))^{\!*}$
  for the dual of the projection zeroing the $V(H)$-columns and
  $\mathrm{sc} = (V(G) \setminus V(H)) \cup \{r\}$ for the surviving body set.
  Then there is a nonzero polynomial $Q$ over the body-normal
  parameters of $G$ such that at every $Q$-non-root seed $q$, the
  $D$-projected surviving-edge panel rows of $\operatorname{ofNormals}(G
  \setminus E(H), \text{ends}, q)$ contain a linearly independent subfamily of
  size at least $D(|\mathrm{sc}|-1) - k'$.
  This is the surviving-leg rank input Case~I's block-triangular
  argument needs when the contraction leg is not yet fully rigid. It is
  the deficiency-aware twin of the rigid contraction-leg transport
  (\cref{lem:claim-6-4}): where the rigid form derives the surviving rank
  from rigidity and reads it off the un-relabeled rows, this transports
  the contraction's possibly-deficient rank across the collapse relabel,
  projecting away the rigid block's columns.
\end{lemma}
\begin{proof}
  \leanok
  \uses{lem:rank-transport-relabel-of-le-finrank, lem:extProj-preserves-rank-of-inter, lem:rank-polynomial-of-le-finrank}
  \Cref{lem:rank-transport-relabel-of-le-finrank} supplies a
  general-position witness placement at which the relabel-leg
  framework's rigidity-row rank equals $N$ exactly, where
  $(N : \mathbb{Z}) = D(|\mathrm{sc}|-1) - k'$ (the contraction's vertex set is
  $\mathrm{sc}$ and its deficiency is $k'$). At that placement, the
  rigidity-free column-deletion injectivity of
  \cref{lem:extProj-preserves-rank-of-inter} extracts, from the rank $N$,
  a subfamily of relabel-leg links whose $D$-projected collapsed rows are
  independent, of size $\ge N$. The per-edge collapse correspondence
  transports that projected independence to the projected rows of the
  uncollapsed framework at the degenerate placement (KT's $p_2$), giving
  a single-placement witness, which \cref{lem:rank-polynomial-of-le-finrank}'s
  projected form then lifts to the $Q$-non-root rank polynomial.
\end{proof}
